
\section{11. Conclusion}

\subsection{11.1 Summary}

We presented \textbf{Native Sparse Attention} (NSA), which combines a fixed 
sliding window with dynamic top-$k$ global routing. The mask is a union of 
the two, applied before softmax. The implementation is fully vectorized, 
parameter-free, and compatible with four pretrained architectures.

\subsection{11.2 Results Recap}

\begin{table}[H]
\centering
\small
\begin{tabular}{lll}
\toprule
\textbf{Model} & \textbf{Acc $\Delta$} & \textbf{Density} \\
\midrule
DistilBERT & -1.2\% & 63\% \\
RoBERTa & -2.1\% & 63\% \\
GPT-2 & -1.9\% & 63\% \\
\bottomrule
\end{tabular}
\end{table}

\subsection{11.3 Key Takeaways}

1. \textbf{Union of local + global} captures both regimes.
2. \textbf{Zero parameters added.}
3. \textbf{Density drops with sequence length.}
4. \textbf{Selection is content-dependent.}
5. \textbf{No time speedup without sparse kernels.}

\subsection{11.4 Future Work}

\subsubsection{Short-Term}
\begin{itemize}
\item Triton kernels for sparse matmul
\item Learned window and top-k per head
\item Long-context evaluation
\end{itemize}

\subsubsection{Long-Term}
\begin{itemize}
\item Theoretical error bounds
\item Cross-layer selection transfer
\item Integration with quantization and MoE
\end{itemize}

\subsection{11.5 Closing Remarks}

NSA provides a simple, effective, and parameter-free sparse attention 
mechanism for long-context inference.

\subsection{11.6 Acknowledgments}

We thank {{ACKNOWLEDGMENTS}}.

\subsection{11.7 Author Contributions}

\begin{itemize}
\item \textbf{{{AUTHOR\_1}}} designed NSA, implemented the code, wrote the paper.
\item \textbf{{{AUTHOR\_2}}} provided feedback, edited.
\end{itemize}

\subsection{11.8 Conflict of Interest}

None.